\chapter{Flexibility Method}
\label{ch:flexmethod}

\section{Complementary strain energy}
\begin{equation}
  U^* = \sum_{i=1}^m\tfrac12\F_i^\tr\mat{r}_i = \sum_{i=1}^m\tfrac12\F_i^\tr\mat{d}_i\F_i = \tfrac12\F'^\tr\mat{d}'\F'
\end{equation}
where $\F_i$ is the independent force vector (we use the notation $\F_f$ in Ch.~\ref{ch:flex}) for the $i$-th member, $\mat{r}_i$ is the associated deformation vector, and $\mat{d}_i$ is the flexibility matrix. The definitions of $\F'$ and $\mat{d}'$ are given as follows:
\begin{equation}
  \F' = \Bmat{\F_1\\ \F_2\\ \vdots\\ \F_m},\quad\text{and}\quad \mat{d}' = \bmat{\mat{d}_1 & & & \\ & \mat{d}_2 & & \\ & & \ddots & \\ & & & \mat{d}_m}
\end{equation}
Note that $\F_i$ may include the reaction(s) of a support; in this case, $\mat{d}_i = \mat{0}$ if the support is rigid.

In the flexibility method, certain redundant forces, with the number being equal to redundancy, need to be chosen from $\F'$ first. Let these redundant forces be denoted by $\F_x$. Group $\F'$ into $\F_x$ and the other forces to yield
\begin{equation}
  \F = \Bmat{\F_0\\ \F_x}
\end{equation}
where $\F_0 = \F'$ excludes $\F_x$.

The complementary strain energy, written in terms of $\F$, takes the form
\begin{equation}
  U^* = \tfrac12\F^\tr\mat{d}\F
\end{equation}
The flexibility matrix $\mat{d}$ corresponding to $\F$ can also be obtained by rearranging $\mat{d}'$. Let $\mat{d}$ be expressed as
\begin{equation}
  \mat{d} = \bmat{\mat{d}_{00} & \mat{d}_{0x}\\ \mat{d}_{x0} & \mat{d}_{xx}}
\end{equation}
Note that:
\begin{enumerate}
  \item $\mat{d}$ is symmetric.
  \item $\mat{d}_{0x} = \mat{d}_{x0} = \mat{0}$ for trusses.
  \item In general, $\mat{d}_{0x} = \mat{d}_{x0}\neq\mat{0}$ for beams.
\end{enumerate}

\section{Equilibrium equations}
On the basis of the equilibrium equations for the primary structure (a statically determinate structure), we can write $\F_0$ in terms of $\Pvec_0$ and $\F_x$ as follows.
\begin{equation}
  \F_0 = \mat{C}_1\Pvec_0 + \mat{C}_2\F_x
\end{equation}
where $\Pvec_0 =$ external load vector, and $\mat{C}_1$ and $\mat{C}_2$ are related to the geometry of structures. The $i$-th column of $\mat{C}_1$ represents the forces $\F_0$ due to $(\Pvec_0)_i = 1$ and the $i$-th column of $\mat{C}_2$ represents the forces $\F_0$ due to $(\F_x)_i = 1$.

With the above equilibrium equations, we can write $\F$ in the following form.
\begin{equation}
  \F = \Bmat{\F_0\\ \F_x} = \bmat{\mat{C}_1 & \mat{C}_2\\ \mat{0} & \mat{I}}\Bmat{\Pvec_0\\ \F_x} = \mat{C}\Pvec
\end{equation}
where
\begin{equation}
  \mat{C} = \bmat{\mat{C}_1 & \mat{C}_2\\ \mat{0} & \mat{I}}\quad\text{and}\quad \Pvec = \Bmat{\Pvec_0\\ \F_x}
\end{equation}

\section{Compatibility equations}
Substituting the above expression for $\F$ into $U^*$ results in
\begin{equation}
  U^* = \tfrac12\F^\tr\mat{d}\F = \tfrac12(\mat{C}\Pvec)^\tr\mat{d}(\mat{C}\Pvec) = \tfrac12\Pvec^\tr\mat{D}\Pvec
\end{equation}
where the structural flexibility matrix $\mat{D}$ is defined as
\begin{equation}
  \mat{D} = \mat{C}^\tr\mat{d}\mat{C} = \bmat{\mat{C}_1^\tr & \mat{0}\\ \mat{C}_2^\tr & \mat{I}}\bmat{\mat{d}_{00} & \mat{d}_{0x}\\ \mat{d}_{x0} & \mat{d}_{xx}}\bmat{\mat{C}_1 & \mat{C}_2\\ \mat{0} & \mat{I}} = \bmat{\mat{D}_{pp} & \mat{D}_{px}\\ \mat{D}_{xp} & \mat{D}_{xx}}
\end{equation}
where
\begin{align}
  \mat{D}_{pp} &= \mat{C}_1^\tr\mat{d}_{00}\mat{C}_1\\
  \mat{D}_{px} &= \mat{C}_1^\tr(\mat{d}_{00}\mat{C}_2 + \mat{d}_{0x}) = \mat{D}_{xp}^\tr\\
  \mat{D}_{xx} &= \mat{C}_2^\tr(\mat{d}_{00}\mat{C}_2 + \mat{d}_{0x}) + \mat{d}_{x0}\mat{C}_2 + \mat{d}_{xx}
\end{align}
Note that $\mat{D}$ is symmetric.

According to Castigliano's second theorem, we obtain the following compatibility equations:
\begin{equation}
  \pd{U^*}{P_{0i}} = \Delta_{0i},\quad i = 1, n
\end{equation}
where $n =$ dimension of $\Pvec_0$.

On the other hand we obtain, according to the theorem of least work, the following compatibility equations:
\begin{equation}
  \pd{U^*}{F_{xi}} = 0,\quad i = 1, r
\end{equation}
where $r =$ the number of redundant forces.

The above two sets of equations can be expressed, in vector form, as
\begin{equation}
  \pd{U^*}{\Pvec} = \Bmat{\bs{\Delta}_0\\ \mat{0}}
\end{equation}
Knowing that $U^* = \frac12\Pvec^\tr\mat{D}\Pvec$, therefore $\partial U^*/\partial\Pvec = \mat{D}\Pvec$, and the above equations become
\begin{equation}
  \mat{D}\Pvec = \Bmat{\bs{\Delta}_0\\ \mat{0}},\quad\text{i.e.}\quad
  \bmat{\mat{D}_{pp} & \mat{D}_{px}\\ \mat{D}_{xp} & \mat{D}_{xx}}\Bmat{\Pvec_0\\ \F_x} = \Bmat{\bs{\Delta}_0\\ \mat{0}}
\end{equation}
Solving for $\F_x$ from the lower part of equations results in
\begin{equation}
  \F_x = -\mat{D}_{xx}^{-1}\mat{D}_{xp}\Pvec_0
\end{equation}
Substituting the above expression for $\F_x$ into the upper part of equations yields
\begin{equation}
  \bs{\Delta}_0 = (\mat{D}_{pp} - \mat{D}_{px}\mat{D}_{xx}^{-1}\mat{D}_{xp})\Pvec_0 = \widehat{\mat{D}}_{pp}\Pvec_0
\end{equation}
where
\begin{equation}
  \widehat{\mat{D}}_{pp} = \mat{D}_{pp} - \mat{D}_{px}\mat{D}_{xx}^{-1}\mat{D}_{xp}
\end{equation}
Inverting the relation, $\bs{\Delta}_0 = \widehat{\mat{D}}_{pp}\Pvec_0$, yields the stiffness equations:
\begin{equation}
  \widehat{\Kmat}_{pp}\bs{\Delta}_0 = \Pvec_0
\end{equation}
where $\widehat{\Kmat}_{pp} = \widehat{\mat{D}}_{pp}^{-1}$. Note that $\widehat{\Kmat}_{pp}$ is the condensed stiffness matrix corresponding to $\Pvec_0$ if the stiffness method is employed.

\begin{example}[Truss]
Fig.~(a) shows a truss structure, and Fig.~(b) shows the primary structure, where $F_{bd}$ is the redundant force. Assuming all bars have the same $EA$, determine $F_{bd}$ and $u_c$.
\begin{fig}
  % Truss (a) and primary structure (b) with redundant F_bd
\begin{tikzpicture}[scale=0.8,font=\small]
  \begin{scope}
    \coordinate (a) at (0,0); \coordinate (b) at (0,4); \coordinate (c) at (4,4); \coordinate (d) at (4,0);
    \draw[member] (a) -- (b) -- (c) -- (d); \draw[member] (a) -- (c); \draw[member] (b) -- (d);
    \pic at (a) {pin}; \pic at (d) {pin}; \foreach \p in {a,b,c,d}{\node[node dot] at (\p) {};}
    \node[left] at (a) {$a$}; \node[above left] at (b) {$b$}; \node[above right] at (c) {$c$}; \node[right] at (d) {$d$};
    \draw[load] (c) -- ++(1.4,0) node[right] {$P$};
    \draw[dim] (-1.0,0) -- (-1.0,4) node[midway,left] {$L$}; \draw[dim] (0,-1.0) -- (4,-1.0) node[midway,fill=white,inner sep=1pt] {$L$};
    \node[below] at (2,-1.5) {(a)};
  \end{scope}
  \begin{scope}[xshift=8cm]
    \coordinate (a) at (0,0); \coordinate (b) at (0,4); \coordinate (c) at (4,4); \coordinate (d) at (4,0);
    \draw[member] (a) -- (b) -- (c) -- (d); \draw[member] (a) -- (c);
    \pic at (a) {pin}; \pic at (d) {pin}; \foreach \p in {a,b,c,d}{\node[node dot] at (\p) {};}
    \node[left] at (a) {$a$}; \node[above left] at (b) {$b$}; \node[above right] at (c) {$c$}; \node[right] at (d) {$d$};
    \draw[load] (c) -- ++(1.4,0) node[right] {$P$};
    \draw[load] (b) -- ++(-45:1.4); \draw[load] (d) -- ++(135:1.4); \node at (1.9,2.1) {$F_{bd}$};
    \draw[dim] (-1.0,0) -- (-1.0,4) node[midway,left] {$L$}; \draw[dim] (0,-1.0) -- (4,-1.0) node[midway,fill=white,inner sep=1pt] {$L$};
    \node[below] at (2,-1.5) {(b) Primary structure};
  \end{scope}
\end{tikzpicture}

\end{fig}
First, find the member forces for the primary structure in terms of $P$ and $F_{bd}$.
\begin{equation*}
  \F_0 = \Bmat{F_{ab}\\ F_{bc}\\ F_{cd}\\ F_{ac}} = \Bmat{0\\ 0\\ -1\\ \sqrt2}P + \Bmat{-1\\ -1\\ -1\\ \sqrt2}\frac{F_{bd}}{\sqrt2} = \mat{C}_1\Pvec_0 + \mat{C}_2\F_x,
\end{equation*}
where $\Pvec_0 = P$ and $\F_x = F_{bd}$.

Therefore,
\begin{equation*}
  \mat{C} = \bmat{\mat{C}_1 & \mat{C}_2\\ \mat{0} & \mat{I}} =
  \begin{bmatrix}
    0 & -\dfrac{1}{\sqrt2}\\[6pt] 0 & -\dfrac{1}{\sqrt2}\\[6pt] -1 & -\dfrac{1}{\sqrt2}\\[6pt] \sqrt2 & 1\\[4pt] 0 & 1
  \end{bmatrix}
\end{equation*}
The compatibility equations are
\begin{equation*}
  \mat{D}\Pvec = \Bmat{\bs{\Delta}_0\\ 0} \Rightarrow \mat{D}\Bmat{P\\ F_{bd}} = \Bmat{u_c\\ 0}
\end{equation*}
where $\mat{D} = \mat{C}^\tr\mat{d}\mat{C} = \dfrac{L}{EA}\begin{bmatrix}1+2\sqrt2 & 2+\dfrac{1}{\sqrt2}\\[6pt] 2+\dfrac{1}{\sqrt2} & \dfrac32+2\sqrt2\end{bmatrix}$ since
\begin{equation*}
  \mat{d} = \begin{bmatrix}
    \dfrac{L}{EA} & & & & \\[4pt] & \dfrac{L}{EA} & & & \\[4pt] & & \dfrac{L}{EA} & & \\[4pt] & & & \dfrac{\sqrt2L}{EA} & \\[4pt] & & & & \dfrac{\sqrt2L}{EA}
  \end{bmatrix}
\end{equation*}
Hence, $F_{bd}$ and $u_c$ can be solved from the compatibility equations; the results are:
\begin{equation*}
  \Bmat{F_{bd}\\ u_c} = \Bmat{\dfrac{-(4+\sqrt2)}{3+4\sqrt2}\\[8pt] \dfrac{L}{EA}\dfrac{10+6\sqrt2}{3+4\sqrt2}}P = \Bmat{-0.625P\\[4pt] 2.135\dfrac{PL}{EA}}
\end{equation*}
The axial forces in the other members:
\begin{equation*}
  \Bmat{F_{ab}\\ F_{bc}\\ F_{cd}\\ F_{ac}} = \mat{C}_1\Pvec_0 + \mat{C}_2\F_x = \Bmat{0.442\\ 0.442\\ -0.558\\ 0.789}P
\end{equation*}
Since only a horizontal load is applied at node $C$, the displacement that we can find corresponds to this load. If the displacements of other points or $v_c$ are needed, we shall add artificial loads first and then let them equal zero.
\end{example}

\begin{example}[Beam]
The beam is subjected to a load $P$ at the central point $b$; determine the vertical displacement of node $b$.
\begin{fig}
  % Propped cantilever with central load P and the element end moments
\begin{tikzpicture}[scale=0.7,font=\small]
  \begin{scope}
    \pic at (0,0) {wallL}; \draw[member] (0,0) -- (6,0); \pic at (6,0) {roller}; \node[node dot] at (6,0) {};
    \node[above right] at (0,0.05) {$a$}; \node[above right] at (3,0.05) {$b$}; \node[above left] at (6,0.05) {$C$};
    \draw[load] (3,1.3) -- (3,0.1) node[pos=0,above] {$P$};
    \draw[dim] (0,-1.0) -- (6,-1.0) node[midway,fill=white,inner sep=1pt] {$L$};
  \end{scope}
  \begin{scope}[xshift=8cm]
    \draw[thinmember] (0,0) -- (3.5,0); \draw[thinmember] (6.5,0) -- (10,0);
    \node[above right] at (0.05,0) {$a$}; \node[above left] at (3.5,0) {$b$}; \node[above right] at (6.55,0) {$b$}; \node[above left] at (10,0) {$C$};
    \draw[->,thick] (-0.15,0.45) arc (110:-110:0.45); \node[above] at (-0.2,0.6) {$M_1^{ab}$};
    \draw[->,thick] (3.65,0.45) arc (70:290:0.45); \node[above] at (3.7,0.6) {$M_2^{ab}$};
    \draw[->,thick] (6.35,0.45) arc (110:-110:0.45); \node[above] at (6.4,0.6) {$M_1^{bc}$};
    \draw[->,thick] (10.15,0.45) arc (70:290:0.45); \node[above] at (10.2,0.6) {$M_2^{bc}$};
  \end{scope}
\end{tikzpicture}

\end{fig}
\emph{Case 1: Take the reaction at the roller as the redundant force.}
\begin{fig}
  % Case 1: reaction at the roller as the redundant force
\begin{tikzpicture}[scale=0.7,font=\small]
  \begin{scope}
    \pic at (0,0) {wallL}; \draw[member] (0,0) -- (6,0);
    \node[above right] at (0,0.05) {$a$}; \node[above right] at (3,0.05) {$b$}; \node[above left] at (6,0.05) {$C$};
    \draw[load] (3,1.3) -- (3,0.1) node[pos=0,above] {$P$};
    \draw[load] (6,-1.2) -- (6,-0.1) node[pos=0,below] {$R$};
  \end{scope}
  \begin{scope}[xshift=8cm]
    \draw[thinmember] (0,0) -- (3.5,0); \draw[thinmember] (6.5,0) -- (10,0);
    \node[above right] at (0.05,0) {$a$}; \node[above left] at (3.5,0) {$b$}; \node[above right] at (6.55,0) {$b$}; \node[above left] at (10,0) {$C$};
    \draw[->,thick] (-0.15,0.45) arc (110:-110:0.45); \node[above] at (-0.2,0.6) {$\frac{PL}{2}-RL$};
    \draw[->,thick] (3.65,0.45) arc (70:290:0.45); \node[above] at (3.7,0.6) {$\frac{RL}{2}$};
    \draw[->,thick] (6.35,0.45) arc (110:-110:0.45); \node[above] at (6.4,0.6) {$-\frac{RL}{2}$};
    \draw[->,thick] (10.15,0.45) arc (70:290:0.45); \node[above] at (10.2,0.6) {$0$};
  \end{scope}
\end{tikzpicture}

\end{fig}
First let $\F = \Bmat{\F_0\\ \F_x} = \Bmat{M_1^{ab}\\ M_2^{ab}\\ M_1^{bc}\\ M_2^{bc}\\ R}$. Then analyze this beam and express the element end moments in terms of $P$ and $R$; namely
\begin{equation*}
  \Bmat{M_1^{ab}\\ M_2^{ab}\\ M_1^{bc}\\ M_2^{bc}} = \Bmat{L/2\\ 0\\ 0\\ 0}P + \Bmat{-L\\ L/2\\ -L/2\\ 0}R = \mat{C}_1\Pvec_0 + \mat{C}_2\F_x,\quad\text{where } \mat{C}_1 = \Bmat{L/2\\ 0\\ 0\\ 0}\text{ and } \mat{C}_2 = \Bmat{-L\\ L/2\\ -L/2\\ 0}.
\end{equation*}
Note that $\Pvec_0 = P$ and $\F_x = R$.

We can obtain the structural flexibility matrix as
\begin{equation*}
  \mat{D} = \mat{C}^\tr\mat{d}\mat{C} = \frac{L^3}{EI}\bmat{1/24 & -5/48\\ -5/48 & 1/3}
\end{equation*}
where
\begin{equation*}
  \mat{d} = \bmat{\mat{d}_{ab} & & \\ & \mat{d}_{bc} & \\ & & 0} = \frac{L'}{6EI}\bmat{2 & -1 & & & \\ -1 & 2 & & & \\ & & 2 & -1 & \\ & & -1 & 2 & \\ & & & & 0}\quad\text{with } L' = \frac L2
\end{equation*}
and
\begin{equation*}
  \mat{C}_{5\times2} = \bmat{\mat{C}_1 & \mat{C}_2\\ 0 & 1}
\end{equation*}
Note that the flexibility coefficient for the roller ($d_{55}$) is 0 since it is rigid. Finally, the system compatibility equations can be written as
\begin{equation*}
  \mat{D}\Bmat{P\\ R} = \Bmat{v_b\\ 0},
\end{equation*}
from which we obtain
\begin{equation*}
  \Bmat{R\\ v_b} = \Bmat{\dfrac{5}{16}P\\[6pt] 9.11\times10^{-3}\dfrac{PL^3}{EI}}
\end{equation*}

\emph{Case 2: Take the moment at the fixed end as the redundant force.}
\begin{fig}
  % Case 2: moment at the fixed end as the redundant force
\begin{tikzpicture}[scale=0.7,font=\small]
  \begin{scope}
    \draw[member] (0,0) -- (6,0); \pic at (0,0) {pin}; \pic at (6,0) {roller}; \node[node dot] at (0,0) {}; \node[node dot] at (6,0) {};
    \node[above right] at (0.1,0.05) {$a$}; \node[above right] at (3,0.05) {$b$}; \node[above left] at (5.9,0.05) {$C$};
    \draw[load] (3,1.3) -- (3,0.1) node[pos=0,above] {$P$};
    \draw[->,thick] (-0.5,0.7) arc (100:260:0.5); \node[above] at (-0.6,0.75) {$R$};
    \draw[dim] (0,-1.0) -- (6,-1.0) node[midway,fill=white,inner sep=1pt] {$L$};
  \end{scope}
  \begin{scope}[xshift=8cm]
    \draw[thinmember] (0,0) -- (3.5,0); \draw[thinmember] (6.5,0) -- (10,0);
    \node[above right] at (0.05,0) {$a$}; \node[above left] at (3.5,0) {$b$}; \node[above right] at (6.55,0) {$b$}; \node[above left] at (10,0) {$C$};
    \draw[->,thick] (-0.15,0.45) arc (110:-110:0.45); \node[above] at (-0.2,0.6) {$R$};
    \draw[->,thick] (3.65,0.45) arc (70:290:0.45); \node[above] at (3.5,0.6) {$\frac{PL}{4}-\frac R2$};
    \draw[->,thick] (6.35,0.45) arc (110:-110:0.45); \node[above] at (6.8,0.6) {$-\frac{PL}{4}+\frac R2$};
    \draw[->,thick] (10.15,0.45) arc (70:290:0.45); \node[above] at (10.2,0.6) {$0$};
  \end{scope}
\end{tikzpicture}

\end{fig}
Following the presented procedure, we will have
\begin{equation*}
  \mat{C}_1 = \Bmat{0\\ \dfrac L4\\[4pt] -\dfrac L4\\[4pt] 0},\quad \mat{C}_2 = \Bmat{1\\ -\dfrac12\\[4pt] \dfrac12\\[4pt] 0}\quad\text{and}\quad \mat{D} = \frac{L^3}{EI}\begin{bmatrix}\dfrac{1}{48} & -\dfrac{1}{16L}\\[8pt] -\dfrac{1}{16L} & \dfrac{1}{3L^2}\end{bmatrix}.
\end{equation*}
Finally, we have
\begin{equation*}
  \Bmat{R\\ v_b} = \Bmat{\dfrac{3}{16}PL\\[6pt] 9.11\times10^{-3}\dfrac{PL^3}{EI}}
\end{equation*}
The above two cases use end moments as the DOFs. When the moment and force of the right end are chosen to be the DOFs, we can obtain the same solutions. For example, when the fixed end moment is chosen as the redundant force (case 2), we have
\begin{equation*}
  \mat{d} = \bmat{\mat{d}_{ab} & & \\ & \mat{d}_{bc} & \\ & & 0},\qquad \mat{d}_{ab} = \mat{d}_{bc} = \frac{L'}{6EI}\bmat{2L'^2 & 3L'\\ 3L' & 6},\ \text{where } L' = \frac L2.
\end{equation*}
\begin{equation*}
  \mat{C}_1 = \Bmat{-\dfrac12\\[4pt] \dfrac L4\\[4pt] \dfrac12\\[4pt] 0},\quad \mat{C}_2 = \Bmat{-\dfrac1L\\[4pt] -\dfrac12\\[4pt] -\dfrac1L\\[4pt] 0},\quad\text{and}\quad \mat{D} = \frac{L^3}{EI}\begin{bmatrix}\dfrac{1}{48} & -\dfrac{1}{16L}\\[8pt] -\dfrac{1}{16L} & \dfrac{1}{3L^2}\end{bmatrix}
\end{equation*}
Note that the same $\mat{D}$ is obtained.
\end{example}

\begin{example}[Beam with an internal spring]
The structure is similar to the previous example except that an internal rotational spring is added at the loaded point.
\begin{fig}
  % Propped cantilever with an internal rotational spring at the loaded point
\begin{tikzpicture}[scale=0.8,font=\small]
  \pic at (0,0) {wallL}; \draw[member] (0,0) -- (6,0); \pic at (6,0) {roller}; \node[node dot] at (6,0) {};
  \draw[thick] (2,0) circle (0.18); \draw[thick] (2,0) circle (0.09);
  \node[above right] at (0,0.05) {$a$}; \node[above right] at (2.1,0.1) {$b$}; \node[above left] at (6,0.05) {$C$};
  \draw[load] (2,0.2) -- (2,1.3) node[above] {$P$};
  \draw[dim] (0,-1.0) -- (2,-1.0) node[midway,fill=white,inner sep=1pt] {$L$}; \draw[dim] (2,-1.0) -- (6,-1.0) node[midway,fill=white,inner sep=1pt] {$2L$};
\end{tikzpicture}

\end{fig}
Take the reaction at the roller as the redundant force.
\begin{fig}
  % Primary structure (redundant R at the roller) and element end moments incl. the spring moment
\begin{tikzpicture}[scale=0.8,font=\small]
  \begin{scope}
    \pic at (0,0) {wallL}; \draw[member] (0,0) -- (6,0);
    \draw[thick] (2,0) circle (0.18); \draw[thick] (2,0) circle (0.09);
    \node[above right] at (0,0.05) {$a$}; \node[above right] at (2.1,0.1) {$b$}; \node[above left] at (6,0.05) {$C$};
    \draw[load] (2,0.2) -- (2,1.3) node[above] {$P$};
    \draw[load] (6,-1.2) -- (6,-0.1) node[pos=0,left] {$R$};
  \end{scope}
  \begin{scope}[xshift=7.5cm]
    \draw[thinmember] (0,0) -- (2.2,0); \draw[thick] (3.1,0) circle (0.16); \draw[thick] (3.1,0) circle (0.08); \draw[thinmember] (4.0,0) -- (8.0,0);
    \node[above right] at (0.05,0) {$a$}; \node[above left] at (2.2,0) {$b$}; \node[above right] at (4.05,0) {$b$}; \node[above left] at (8.0,0) {$C$};
    \draw[->,thick] (-0.15,0.45) arc (110:-110:0.45); \node[above] at (-0.2,0.6) {$M_1^{ab}$};
    \draw[->,thick] (2.35,0.45) arc (70:290:0.45); \node[above] at (2.4,0.6) {$M_2^{ab}$};
    \draw[->,thick] (3.6,0.45) arc (70:290:0.45); \node[above] at (3.4,0.6) {$M_s$};
    \draw[->,thick] (3.85,0.45) arc (110:-110:0.45); \node[above] at (4.4,0.6) {$M_1^{bc}$};
    \draw[->,thick] (8.15,0.45) arc (70:290:0.45); \node[above] at (8.2,0.6) {$M_2^{bc}$};
  \end{scope}
\end{tikzpicture}

\end{fig}
\begin{equation*}
  \F = \Bmat{\F_0\\ \F_x} = \Bmat{M_1^{ab}\\ M_2^{ab}\\ M_1^{bc}\\ M_2^{bc}\\ M_s\\ R},\quad\text{and}\quad
  \mat{d} = \bmat{\mat{d}_{ab} & & & \\ & \mat{d}_{bc} & & \\ & & \dfrac{1}{k_\theta} & \\ & & & 0}
\end{equation*}
Using the same procedure, we will obtain $v_b = \dfrac19\dfrac{PL^3}{EI}\dfrac{11\alpha+12}{9\alpha+4}$. Try it.
\end{example}

\begin{example}[2D frame]
\leavevmode

\begin{minipage}{0.6\linewidth}
Even for 2D frames, the same procedure can be used. The DOFs for one member are shown in the figure, and the flexibility matrix is
\begin{equation*}
  \mat{d} = \begin{bmatrix}\dfrac{L}{EA} & & \\[6pt] & \dfrac{L}{3EI} & -\dfrac{L}{6EI}\\[8pt] & -\dfrac{L}{6EI} & \dfrac{L}{3EI}\end{bmatrix}
\end{equation*}
\end{minipage}\hfill
\begin{minipage}{0.35\linewidth}\centering% 2D frame: the three member DOFs (axial + two end rotations) of the left column
\begin{tikzpicture}[scale=0.8,font=\small]
  \draw[member] (0,0) -- (0,4) -- (4,4) -- (4,0);
  \pic at (0,0) {pin}; \pic at (4,0) {fixed}; \node[node dot] at (0,0) {};
  \draw[load] (0,4.1) -- (0,5.2);
  \draw[->,thick] (-0.5,4.4) arc (150:-30:0.5);
  \draw[->,thick] (0.5,0.4) arc (30:210:0.5);
\end{tikzpicture}
\end{minipage}

This problem was solved by the stiffness method in Sec.~\ref{sec:axial-portal}.
\begin{fig}
  % Portal frame solved by the flexibility method: redundants at d
\begin{tikzpicture}[scale=0.8,font=\small]
  \draw[member] (0,0) -- (0,4) -- (4,4) -- (4,0);
  \pic at (0,0) {fixed};
  \node[left] at (0,0.3) {$a$}; \node[above right] at (0,4) {$b$}; \node[above left] at (4,4) {$c$}; \node[above left] at (4,0.3) {$d$};
  \draw[load] (-1.4,4) -- (-0.1,4) node[pos=0,below] {$P$};
  \draw[load] (4.1,0) -- (5.4,0) node[right] {$F_{x1}$};
  \draw[load] (4,-1.3) -- (4,-0.1) node[pos=0,right] {$F_{x2}$};
  \draw[->,thick] (4.6,-0.4) arc (-40:200:0.6); \node[left] at (3.4,-0.4) {$F_{x3}$};
\end{tikzpicture}

\end{fig}
\begin{equation*}
  \mat{C} = \bigl[(\mat{C}_1)_{9\times1}, (\mat{C}_2)_{9\times3}\bigr] =
  \begin{bmatrix}
    0 & 0 & 1 & 0\\ L & 0 & -L & -1\\ 0 & L & L & 1\\ 0 & 1 & 0 & 0\\ 0 & -L & -L & -1\\ 0 & L & 0 & 1\\ 0 & 0 & -1 & 0\\ 0 & -L & 0 & -1\\ 0 & 0 & 0 & 1
  \end{bmatrix}
\end{equation*}
\begin{equation*}
  \F_0 = \Bmat{\F^{ab}\\ \F^{bc}\\ \F^{cd}};\quad\text{member forces of each member: } \F = \Bmat{F\\ M_1\\ M_2}
\end{equation*}
{\small
\begin{equation*}
  \mat{d} = \bmat{\mat{d}_{ab} & \mat{0} & \mat{0}\\ \mat{0} & \mat{d}_{bc} & \mat{0}\\ \mat{0} & \mat{0} & \mat{d}_{cd}} =
  \begin{bmatrix}
    \frac{L}{EA} & 0 & 0 & 0 & 0 & 0 & 0 & 0 & 0\\
    0 & \frac{L}{3EI} & \frac{-L}{6EI} & 0 & 0 & 0 & 0 & 0 & 0\\
    0 & \frac{-L}{6EI} & \frac{L}{3EI} & 0 & 0 & 0 & 0 & 0 & 0\\
    0 & 0 & 0 & \frac{L}{EA} & 0 & 0 & 0 & 0 & 0\\
    0 & 0 & 0 & 0 & \frac{L}{3EI} & \frac{-L}{6EI} & 0 & 0 & 0\\
    0 & 0 & 0 & 0 & \frac{-L}{6EI} & \frac{L}{3EI} & 0 & 0 & 0\\
    0 & 0 & 0 & 0 & 0 & 0 & \frac{L}{EA} & 0 & 0\\
    0 & 0 & 0 & 0 & 0 & 0 & 0 & \frac{L}{3EI} & \frac{-L}{6EI}\\
    0 & 0 & 0 & 0 & 0 & 0 & 0 & \frac{-L}{6EI} & \frac{L}{3EI}
  \end{bmatrix}
\end{equation*}
\begin{equation*}
  \mat{D} = \mat{C}^\tr\mat{d}\mat{C} =
  \begin{bmatrix}
    \frac13\frac{L^3}{EI} & -\frac16\frac{L^3}{EI} & -\frac12\frac{L^3}{EI} & -\frac12\frac{L^2}{EI}\\[4pt]
    -\frac16\frac{L^3}{EI} & \frac53\frac{L^3}{EI}+\frac{L}{EA} & \frac{L^3}{EI} & 2\frac{L^2}{EI}\\[4pt]
    -\frac12\frac{L^3}{EI} & \frac{L^3}{EI} & 2\frac{L}{EA}+\frac43\frac{L^3}{EI} & \frac32\frac{L^2}{EI}\\[4pt]
    -\frac12\frac{L^2}{EI} & 2\frac{L^2}{EI} & \frac32\frac{L^2}{EI} & 3\frac{L}{EI}
  \end{bmatrix};\qquad
  \Pvec = \Bmat{P\\ F_{x1}\\ F_{x2}\\ F_{x3}}
\end{equation*}}
System flexibility equations:
\begin{equation*}
  \mat{D}\Pvec = \Bmat{u_b\\ 0\\ 0\\ 0}\quad\rightarrow\quad\text{the last three equations:}
\end{equation*}
\begin{equation*}
  \begin{bmatrix}
    \frac53\frac{L^3}{EI}+\frac{L}{EA} & \frac{L^3}{EI} & 2\frac{L^2}{EI}\\[4pt]
    \frac{L^3}{EI} & 2\frac{L}{EA}+\frac43\frac{L^3}{EI} & \frac32\frac{L^2}{EI}\\[4pt]
    2\frac{L^2}{EI} & \frac32\frac{L^2}{EI} & 3\frac{L}{EI}
  \end{bmatrix}\Bmat{F_{x1}\\ F_{x2}\\ F_{x3}} = \Bmat{\frac16\frac{L^3}{EI}\\[4pt] \frac12\frac{L^3}{EI}\\[4pt] \frac12\frac{L^2}{EI}}P
\end{equation*}
\begin{equation*}
  \rightarrow\quad \Bmat{F_{x1}\\ F_{x2}\\ F_{x3}} = \Bmat{-\frac12L^2\frac{A}{L^2A+3I}\\[6pt] 3L^2\frac{A}{24I+7L^2A}\\[6pt] L\frac{11L^2AI+2L^4A^2+12I^2}{(24I+7L^2A)(L^2A+3I)}}P
\end{equation*}
$\rightarrow$ Substituting them into the first equation obtains:
\begin{equation*}
  u_b = \frac{1}{12}L^3\frac{84L^2AI+216I^2+5L^4A^2}{EI(45L^2AI+72I^2+7L^4A^2)}
\end{equation*}
The above solution is the same as that obtained in Sec.~\ref{sec:axial-portal}.

\topic{Neglect of axial deformation:}
\begin{equation*}
  \frac{L}{EA} = 0\ \rightarrow\ \mat{d} =
  \begin{bmatrix}
    \frac{L}{3EI} & \frac{-L}{6EI} & 0 & 0 & 0 & 0\\
    \frac{-L}{6EI} & \frac{L}{3EI} & 0 & 0 & 0 & 0\\
    0 & 0 & \frac{L}{3EI} & \frac{-L}{6EI} & 0 & 0\\
    0 & 0 & \frac{-L}{6EI} & \frac{L}{3EI} & 0 & 0\\
    0 & 0 & 0 & 0 & \frac{L}{3EI} & \frac{-L}{6EI}\\
    0 & 0 & 0 & 0 & \frac{-L}{6EI} & \frac{L}{3EI}
  \end{bmatrix};
\end{equation*}
\begin{equation*}
  \mat{C} = \bigl[(\mat{C}_1)_{6\times1}, (\mat{C}_2)_{6\times3}\bigr] =
  \begin{bmatrix}
    L & 0 & -L & -1\\ 0 & L & L & 1\\ 0 & -L & -L & -1\\ 0 & L & 0 & 1\\ 0 & -L & 0 & -1\\ 0 & 0 & 0 & 1
  \end{bmatrix}
\end{equation*}
\begin{equation*}
  \mat{D} = \begin{bmatrix}
    \frac13\frac{L^3}{EI} & -\frac16\frac{L^3}{EI} & -\frac12\frac{L^3}{EI} & -\frac12\frac{L^2}{EI}\\[4pt]
    -\frac16\frac{L^3}{EI} & \frac53\frac{L^3}{EI} & \frac{L^3}{EI} & 2\frac{L^2}{EI}\\[4pt]
    -\frac12\frac{L^3}{EI} & \frac{L^3}{EI} & \frac43\frac{L^3}{EI} & \frac32\frac{L^2}{EI}\\[4pt]
    -\frac12\frac{L^2}{EI} & 2\frac{L^2}{EI} & \frac32\frac{L^2}{EI} & 3\frac{L}{EI}
  \end{bmatrix}
\end{equation*}
The last three equations:
\begin{equation*}
  \begin{bmatrix}
    \frac53\frac{L^3}{EI} & \frac{L^3}{EI} & 2\frac{L^2}{EI}\\[4pt]
    \frac{L^3}{EI} & \frac43\frac{L^3}{EI} & \frac32\frac{L^2}{EI}\\[4pt]
    2\frac{L^2}{EI} & \frac32\frac{L^2}{EI} & 3\frac{L}{EI}
  \end{bmatrix}\Bmat{F_{x1}\\ F_{x2}\\ F_{x3}} = \Bmat{\frac16\frac{L^3}{EI}\\[4pt] \frac12\frac{L^3}{EI}\\[4pt] \frac12\frac{L^2}{EI}}P
  \quad\rightarrow\quad \Bmat{F_{x1}\\ F_{x2}\\ F_{x3}} = \Bmat{-\frac12\\[4pt] \frac37\\[4pt] \frac27L}P
\end{equation*}
$\rightarrow$ Substituting them into the first equation obtains: $u_b = \dfrac{5}{84}\dfrac{PL^3}{EI}$.

This value is the same as that given in Sec.~\ref{sec:axial-portal}.
\end{example}
